\documentclass[12pt]{article}
\usepackage{amsmath,amssymb,amsthm}

\newtheorem{theorem}{Theorem}[section]
\newtheorem{lemma}[theorem]{Lemma}
\newtheorem{remark}[theorem]{Remark}

\begin{document}

\section*{Appendix: Admissible-Step Foundations}

\subsection*{Preface: Assumptions and Notation}
This appendix formalizes foundational pullbacks required to evaluate the existence of admissible 3-connectivity-preserving reduction steps. We rely strictly on the definitions from the main manuscript:
\begin{itemize}
    \item $G \in \mathcal{Q}$ is a finite, simple, planar, cubic, bipartite, and strictly 3-connected graph.
    \item $N_{\mathrm{base}} = 12$ is the minimum graph size evaluated ($|V(G)| > 12$).
    \item A \textbf{certified occurrence} defines a local planar patch $H \subset G$ separated by exactly 4 distinct external boundary terminals $B \subset V(G) \setminus V(H)$. Thus, exactly 4 structural edges connect $H$ to $G \setminus H$.
    \item A \textbf{reduction gadget} $H'$ formally replaces $H$. For the $C_2$ configuration, $H'$ consists inherently of exactly two adjacent bridging vertices, $x$ and $y$. The gadget specifies that $x$ uniquely connects to a 2-terminal subset $B_x \subset B$, and $y$ uniquely connects to the remaining 2-terminal subset $B_y \subset B$. The vertices $x$ and $y$ evaluate identically as degree 3 inside the resulting modified graph $G' = (G \setminus H) \cup H'$.
    \item A reduction step evaluates as \textbf{admissible} if applying $H \to H'$ strictly preserves 3-connectivity, successfully returning $G' \in \mathcal{Q}$.
\end{itemize}

\subsection*{Component Linkage and Pullback Mechanism}

\begin{lemma}[Connectedness of $G \setminus H$]
\label{lem:isolated_neighborhood}
Let $G \in \mathcal{Q}$ contain a certified occurrence $H$. Then the induced subgraph
$G[\,V(G)\setminus V(H)\,]$ is connected.
\end{lemma}
\begin{proof}
Assume for contradiction that $G[\,V(G)\setminus V(H)\,]$ has at least two components
$C_1,\dots,C_k$ with $k\ge 2$.
For each $i$, every edge leaving $C_i$ must go to $V(H)$, so $E(C_i,V(H))$ is an edge-cut.
Since $G$ is 3-connected, $\lambda(G)\ge \kappa(G)\ge 3$, hence $|E(C_i,V(H))|\ge 3$ for all $i$.
Therefore
\[
|E(V(H),V(G)\setminus V(H))|
=\sum_{i=1}^k |E(C_i,V(H))|
\ge 3k \ge 6,
\]
contradicting the certified-occurrence fact that exactly $4$ edges join $H$ to $G\setminus H$.
Hence $G[\,V(G)\setminus V(H)\,]$ is connected.
\end{proof}

\begin{lemma}[2-Vertex Cut Pullback for $C_2$]
\label{lem:pullback_cut_c2}
Assume $G\in\mathcal{Q}$ contains a certified $C_2$ occurrence $H$ with boundary
$B=B_x\dot\cup B_y$ (each of size $2$), and let $G'=(G\setminus H)\cup H'$ be the reduced graph,
where $H'$ has two adjacent gadget vertices $x,y$ with $N_{G'}(x)\setminus\{y\}=B_x$ and
$N_{G'}(y)\setminus\{x\}=B_y$.

If $G'$ has a 2-vertex cut $A=\{p,q\}$, then either:
\begin{enumerate}
\item[(i)] $A\subseteq V(G)\setminus V(H)$ and $A$ is a 2-vertex cut of $G$; or
\item[(ii)] $A=\{x,q\}$ for some $q\in V(G)\setminus V(H)$ and $G-(\{q\}\cup B_x)$ is disconnected; or
\item[(iii)] $A=\{y,q\}$ for some $q\in V(G)\setminus V(H)$ and $G-(\{q\}\cup B_y)$ is disconnected.
\end{enumerate}
In cases (ii) and (iii), there is in particular a nonempty vertex set
$K\subseteq V(G)\setminus (V(H)\cup\{q\}\cup B_x)$ (resp.\ $B_y$) that is a union of components of
$G-(\{q\}\cup B_x)$ (resp.\ $G-(\{q\}\cup B_y)$).
\end{lemma}

\begin{proof}
Let $A=\{p,q\}$ be a 2-vertex cut in $G'$, and let $C$ be a component of $G'-A$.

\smallskip
\noindent\textbf{Case 1: $p,q\notin\{x,y\}$.}
Since $xy\in E(G')$, the vertices $x$ and $y$ lie in the same component of $G'-A$; call it $C_0$.
Let $C_1$ be any other component of $G'-A$ (it exists because $A$ is a cut).

First note that $C_1\subseteq V(G)\setminus V(H)$: indeed, the only vertices of $G'$ not in
$V(G)\setminus V(H)$ are $x$ and $y$, and both lie in $C_0$.

Next observe that $B\subseteq V(C_0)\cup A$.
If $b\in B\setminus A$, then $b$ is adjacent in $G'$ to either $x$ or $y$, which lies in $C_0$,
hence $b\in C_0$ as well.

Now consider the original graph $G$. The subgraph $G[\,V(G)\setminus V(H)\,]$ is identical in $G$
and $G'$ (only the interior of $H$ is replaced), so in particular $C_1$ is still a connected subgraph
of $G$ on the same vertex set. Any edge of $G$ leaving $C_1$ either:
(a) stays inside $V(G)\setminus V(H)$, or (b) goes from $V(G)\setminus V(H)$ into $V(H)$ through a
boundary terminal in $B$.
But (a) cannot happen without creating an edge between distinct components of $G'-A$ inside
$V(G)\setminus V(H)$, contradicting that $C_1$ is a component of $G'-A$.
And (b) cannot happen because $B\subseteq V(C_0)\cup A$, so no boundary terminal lies in $C_1$.
Hence every path in $G$ from $C_1$ to $V(G)\setminus C_1$ must meet $A$, i.e.\ $A$ is a 2-vertex cut in $G$.
This proves (i).

\smallskip
\noindent\textbf{Case 2: $\{p,q\}=\{x,y\}$.}
Then $G'-\{x,y\}=G[\,V(G)\setminus V(H)\,]$ (same vertex set and same edges among those vertices).
By Lemma~\ref{lem:isolated_neighborhood}, this graph is connected, so $\{x,y\}$ cannot be a 2-vertex cut.
Contradiction.

\smallskip
\noindent\textbf{Case 3: Exactly one of $p,q$ lies in $\{x,y\}$.}
By symmetry, assume $p=x$ and $q\notin\{x,y\}$.
Let $C_0$ be the component of $G'-\{x,q\}$ that contains $y$.
Since $A$ is a cut, there exists some other component $C_1$ of $G'-\{x,q\}$, disjoint from $C_0$.

We claim that $G-(\{q\}\cup B_x)$ is disconnected.
Define
\[
K := V(C_1)\setminus B_x.
\]
First, $K\neq\emptyset$: if $V(C_1)\cap B_x=\emptyset$ then $K=V(C_1)\neq\emptyset$.
Otherwise pick $b\in V(C_1)\cap B_x$. In $G'$, the vertex $b$ has neighbors $\{x\}\cup\{b_1,b_2\}$
where $b_1,b_2\in V(G)\setminus V(H)$ are its two non-gadget neighbors (they exist because $b$ is cubic
and simple). Since $x$ is deleted and $b\notin\{q\}$, both $b_1,b_2$ remain in $G'-\{x,q\}$ and must lie
in the same component as $b$, hence in $C_1$. At least one of $b_1,b_2$ is not in $B_x$ (because $B_x$
has size $2$ and the graph is simple), so $K$ contains such a neighbor and is nonempty.

Now let $v\in K$ and suppose $v$ has a neighbor $w$ in $G$ with
$w\notin K\cup\{q\}\cup B_x$.
Because $v\notin V(H)$ and $w\notin V(H)$ (as $w\notin B_x$ and edges from $V(G)\setminus V(H)$ into $V(H)$
must go via $B$), the edge $vw$ is also an edge of $G'$ and persists in $G'-\{x,q\}$.
Thus $w$ lies in the same component of $G'-\{x,q\}$ as $v$, namely $C_1$.
Moreover $w\notin B_x$ by assumption, so $w\in K$.
This contradicts $w\notin K\cup\{q\}\cup B_x$.
Therefore, no such neighbor $w$ exists, meaning there are \emph{no} edges in $G$ from $K$ to
$V(G)\setminus (K\cup\{q\}\cup B_x)$.

Equivalently, $K$ is a union of components of $G-(\{q\}\cup B_x)$, and $G-(\{q\}\cup B_x)$ is disconnected.
This is exactly (ii). The case $p=y$ is symmetric and yields (iii).
\end{proof}
\end{document}
